%! Unit 5 — Logarithmic Functions — Unit Cover Page
\documentclass[10pt]{article}
\usepackage{precalculus-article}
\usepackage{precalculus-boxes}
\usepackage{ltablex}
\keepXColumns

\begin{document}
\thispagestyle{empty}

% Full-bleed navy banner
\begin{tikzpicture}[remember picture, overlay]
  \fill[navy] (current page.north west)
    rectangle ([yshift=-1.15in]current page.north east);
\end{tikzpicture}
\vspace{-1.05in}

\begin{center}
  {\color{white}\textbf{\LARGE Precalculus}}\\[6pt]
  {\color{skymid}\Large\textbf{Unit 5: Logarithmic Functions}}\\[4pt]
  {\color{white}\small 8 Lessons + Unit Test \quad|\quad Class Periods: 18--22}
\end{center}

\vspace{0.26in}

\begin{tcolorbox}[colback=goldbg, colframe=goldacc, boxrule=0.8pt, arc=2mm,
    left=3mm, right=3mm, top=2mm, bottom=2mm]
\small\textbf{Unit Essential Questions:}
\begin{itemize}[leftmargin=1.4em, itemsep=2pt, topsep=2pt]
  \item Why do logarithms undo exponentials, and how does that inverse relationship let us solve equations?
  \item How do logarithmic scales make multiplicative change readable?
\end{itemize}
\end{tcolorbox}

\vspace{0.14in}

\begin{tocbox}
\small
\renewcommand{\arraystretch}{1.45}
\begin{tabularx}{\linewidth}{@{} c >{\bfseries}X l @{}}
\rowcolor{sky}
\textbf{\#} & \textbf{Lesson Title} & \textbf{Content Source} \\
\midrule
5.1 & Logarithmic Expressions & CED 2.9 \\
5.2 & Inverses of Exponential Functions & CED 2.10 \\
5.3 & Logarithmic Functions & CED 2.11 \\
5.4 & Logarithmic Function Manipulation & CED 2.12 \\
5.5 & Exponential Equations and Inequalities & Split of CED 2.13 \\
5.6 & Logarithmic Equations and Inequalities & CED 2.13 \\
5.7 & Logarithmic Function Context and Data Modeling & CED 2.14 \\
5.8 & Semi-log Plots & CED 2.15 \\
\midrule
\rowcolor{sky}
     & \textbf{Sample Test} &      \\
\end{tabularx}
\end{tocbox}

\vspace{0.14in}

\begin{skillbox}[Mathematical Practices --- Unit 5 Focus]{goldbox}
\renewcommand{\arraystretch}{1.4}
\begin{tabularx}{\linewidth}{@{} >{\bfseries\color{navy}}p{0.12\linewidth} X @{}}
MP 1.A & Solve equations and inequalities analytically, with and without technology. \\
MP 1.B & Express functions, equations, or expressions in analytically equivalent forms. \\
MP 1.C & Construct new functions using transformations, compositions, inverses, or regressions. \\
MP 2.B & Construct equivalent graphical, numerical, analytical, and verbal representations. \\
MP 3.B & Apply numerical results in a given mathematical or applied context. \\
\end{tabularx}
\end{skillbox}

\end{document}
